% Separate proof increment. All supplied inputs and the checkpoint remain unchanged.
% Unique label prefix: inc:AN02:q2:v1:
\documentclass[11pt]{article}
\usepackage[T1]{fontenc}\usepackage[utf8]{inputenc}\usepackage{lmodern}
\usepackage[margin=0.9in]{geometry}
\usepackage{amsmath,amssymb,amsthm,mathtools,microtype,booktabs,xurl}
\usepackage[hidelinks]{hyperref}\usepackage{fancyhdr}
\pagestyle{fancy}\fancyhf{}\fancyhead[L]{\small AN02: Q2 profile and localized tracking}
\fancyhead[R]{\small Review increment v1}\fancyfoot[C]{\thepage}
\setlength{\headheight}{14pt}\setlength{\emergencystretch}{3em}\allowdisplaybreaks[2]
\numberwithin{equation}{section}
\newtheorem{theorem}{Theorem}[section]\newtheorem{lemma}[theorem]{Lemma}
\newtheorem{proposition}[theorem]{Proposition}\newtheorem{corollary}[theorem]{Corollary}
\theoremstyle{remark}\newtheorem{remark}[theorem]{Remark}
\newcommand{\dd}{\,d}\newcommand{\E}{\mathbb E}\newcommand{\R}{\mathbb R}
\title{Q2 profiles, analytical transit clocks, and localized tracking\\
\large Exact target-only results and an original-flow comparison with an explicit defect}
\author{Analytical increment for independent review}\date{Version 1 --- October 3, 2026}
\hypersetup{pdftitle={Q2 profiles, analytical transit clocks, and localized tracking},pdfauthor={Analytical increment for independent review}}
\begin{document}\maketitle
\begin{abstract}
The positive-noise target-only system admits an exact scalar quadrature and
weighted pushforward formula on each nonstationary angular interval. We
distinguish the Cauchy-squared coordinate-energy profile from total radial
mass and from the angular chart coordinate. For a fixed Q2 label we prove
the weak-axis crossing time, including its finite renormalized constant,
and the leading logarithmic times to a fixed Q1 angle and to the strong
noise-scaled chart. These are target-only statements, not retention or
absorption theorems for the nonlinear population. Separately, an exact
original-flow comparison isolates the principal strong output
$\mathfrak M X_1e_1$. For a left-facing input, that output contributes only
$O(q\mathfrak M)$ to the relevant coordinate and alignment errors. All
remaining output is retained in an explicit $L^2$ defect. The initialized
bound on this defect through strong learning, and the subsequent transit
and return, are still open.
\end{abstract}

\section{Setting and source contract}\label{inc:AN02:q2:v1:setting}
Use normalized distributions
\[
 P_1=\mathcal N(e_1,q^2I_2),\quad
 P_2=\mathcal N(\rho e_2,q^2I_2),\quad P=(P_1+P_2)/2,
 \qquad 0<\rho<1,\quad\lambda=\rho^2.
\]
Write $\phi,\Phi$ for the standard normal density and distribution function,
$K(z)=\phi(z)+z\Phi(z)$, and $a_* =1/2+q^2$.
Initial labels have law $\dd\alpha_0/(2\pi)$ and
$U_0=W_0=\sqrt{S_0}(\cos\alpha_0,\sin\alpha_0)$.
The target-only comparison is
\begin{equation}\label{inc:AN02:q2:v1:target}
 \widehat U'=\E[X(\widehat U^TX)_+],\qquad
 \widehat U=\widehat W=\sqrt{\widehat m}\,a_{\widehat\theta}.
\end{equation}
This is not the original residual equation. Its exact finite-noise
coordinate formula, from Gaussian integration, is
\begin{align}
 \widehat U_j'&=\frac{\mu_j q\sqrt{\widehat m}}2
           K(\mu_j\widehat a_j/q)+d_q\widehat U_j,
           \qquad (\mu_1,\mu_2)=(1,\rho),\nonumber\\
 d_q&=\frac{q^2}{2}[\Phi(\widehat a_1/q)+\Phi(\rho\widehat a_2/q)].
 \label{inc:AN02:q2:v1:coordinates}
\end{align}
The exact original-flow matrices are
\[
 A_\theta=\E[XX^T\mathbf1_{a_\theta^TX>0}],\quad
 C_f(\theta)=\E[f(X)X^T\mathbf1_{a_\theta^TX>0}],\quad
 R_f=A_\theta-C_f.
\]
The foundations and checkpoint give
$0\preceq A_\theta\preceq a_*I$, the balanced characteristic equations,
and the target-only angular derivative cancellations. These equations
are restated below when used. No saved trajectory, numerical fit, or
repository data are a proof premise.

\section{Exact scalar quadrature and density}\label{inc:AN02:q2:v1:quadrature}
Put $s_1=\cos\theta/q$, $s_2=\rho\sin\theta/q$. Equation
\eqref{inc:AN02:q2:v1:coordinates} gives
\begin{align}
 v_q(\theta)&=\frac q2\left[-\sin\theta K(\cos\theta/q)
                 +\rho\cos\theta K(\rho\sin\theta/q)\right],
 \label{inc:AN02:q2:v1:angular}\\
 g_q(\theta)&=q^2\sum_{j=1}^2[s_jK(s_j)+\Phi(s_j)],
 \label{inc:AN02:q2:v1:radial}\\
 \widehat\theta'&=v_q(\widehat\theta),\qquad
 (\log\widehat m)'=g_q(\widehat\theta).\nonumber
\end{align}
Here $g_q$ denotes the full logarithmic mass rate, not half that rate.

\begin{proposition}[First integral and exact weighted pushforward]
\label{inc:AN02:q2:v1:firstintegral}
On an open angular interval $I$ where $v_q$ has no zero, choose a primitive
$G_q$ with $G_q'=g_q/v_q$. Along any characteristic remaining in $I$,
\[
 \log\widehat m(t)-G_q(\widehat\theta(t))
                 =\log S_0-G_q(\alpha_0).
\]
If $\varphi_t$ is the scalar angular flow, its Jacobian is
\[
 \partial_{\alpha_0}\varphi_t(\alpha_0)
           =\frac{v_q(\varphi_t(\alpha_0))}{v_q(\alpha_0)}>0.
\]
Consequently the contribution of such initial labels to the radial
angle measure has density
\begin{equation}\label{inc:AN02:q2:v1:density}
 \frac{\dd\widehat\nu_t}{\dd\theta}
 =\frac{S_0}{2\pi}
 e^{G_q(\theta)-G_q(\alpha_0)}
 \frac{v_q(\alpha_0)}{v_q(\theta)},\qquad
 \alpha_0=\varphi_{-t}(\theta).
\end{equation}
The coordinate-energy density for $\widehat U_2^2$ is
$\sin^2\theta$ times this density. An initial-label restriction adds
its indicator under the inverse flow.
\end{proposition}
\begin{proof}
Differentiate the proposed first integral. For the Jacobian, both the
variational derivative and $v_q(\varphi_t)/v_q(\alpha_0)$ solve
$J'=v_q'(\varphi_t)J$, $J(0)=1$. Smoothness and uniqueness give the
formula. The change of variables from initial labels then gives
\eqref{inc:AN02:q2:v1:density}. All divisions take place on $I$;
no division by a stationary angular speed is permitted.
\end{proof}

\section{What the Cauchy-squared profile measures}
\label{inc:AN02:q2:v1:profile}
This section defines a different comparison: the \emph{zero-noise}
target-only flow on Q2. The positive parameter $q$ below sets the
coordinate scale; it is not noise in this comparison's vector field.
For $\alpha_0\in(\pi/2,\pi)$ the exact solution is
\[
 \overline U_1=\sqrt{S_0}\cos\alpha_0<0,\qquad
 \overline U_2=\sqrt{S_0}\sin\alpha_0e^{\lambda t/2}>0.
\]
Define
\[
 b_t=e^{-\lambda t/2}/q,\qquad
 \upsilon=\frac{\overline U_1}{q\overline U_2}<0,
 \qquad u=\frac{\pi/2-\arg\overline U}{q}.
\]
Thus $\upsilon=\tan(qu)/q$, not exactly $u$.

\begin{proposition}[Coordinate-energy and radial profiles]
\label{inc:AN02:q2:v1:cauchy}
The Q2 coordinate-energy measure
$\overline U_2^2\dd\alpha_0/(2\pi)$ has exact density
\begin{equation}\label{inc:AN02:q2:v1:cauchydensity}
 \dd\nu_{2,t}(\upsilon)
 =\frac{S_0e^{\lambda t}}{2\pi b_t}
      (1+\upsilon^2/b_t^2)^{-2}\dd\upsilon,
      \qquad \upsilon<0.
\end{equation}
Its total mass is $S_0e^{\lambda t}/8$. The total radial measure is
\begin{equation}\label{inc:AN02:q2:v1:radialdensity}
 \dd\nu_{m,t}=(1+q^2\upsilon^2)\dd\nu_{2,t}.
\end{equation}
In the angular coordinate $-\pi/(2q)<u<0$, the radial density is
\begin{equation}\label{inc:AN02:q2:v1:angularprofile}
 \frac{\dd\nu_{m,t}}{\dd u}
 =\frac{S_0e^{\lambda t}}{2\pi b_t}
 \left(1+\frac{\tan^2(qu)}{q^2b_t^2}\right)^{-2}\sec^4(qu).
\end{equation}
The coordinate-energy angular density instead has $\sec^2(qu)$.
At $t=\log(1/S_0)+c$, $b_t=(S_0^{\lambda/2}/q)e^{-\lambda c/2}$
and the prefactor $S_0e^{\lambda t}=S_0^{1-\lambda}e^{\lambda c}$.
\end{proposition}
\begin{proof}
Write $r_0=-\cot\alpha_0>0$, so
$\alpha_0=\pi/2+\arctan r_0$, $\sin^2\alpha_0=(1+r_0^2)^{-1}$,
and $\dd\alpha_0=\dd r_0/(1+r_0^2)$. Since
$\upsilon=-b_t r_0$, the pushforward of
$S_0e^{\lambda t}\sin^2\alpha_0\dd\alpha_0/(2\pi)$ is
\eqref{inc:AN02:q2:v1:cauchydensity}. Its integral uses
$\int_0^\infty(1+r^2)^{-2}\dd r=\pi/4$.
The identity $\overline m=\overline U_2^2(1+q^2\upsilon^2)$
gives \eqref{inc:AN02:q2:v1:radialdensity}, and
$\dd\upsilon/\dd u=\sec^2(qu)$ gives
\eqref{inc:AN02:q2:v1:angularprofile}.
\end{proof}

\paragraph{Finite-noise qualification.}
The Cauchy-squared law is exact for coordinate energy in a tangent
coordinate in this specified zero-noise comparison. For the finite-noise
target-only flow, Proposition \ref{inc:AN02:q2:v1:firstintegral} is the
exact profile. A fixed negative scaled offset still has a nonzero
cross-gate correction $K(u)$ as $q\to0$. Consequently the elementary
profile is not an exact finite-noise law on every $u\le-u_0$ with fixed
$u_0>0$. Both the gate approximation and the angular/radial factors need
control before transferring a density statement.

\begin{proposition}[A quantitative outer Q2 comparison]
\label{inc:AN02:q2:v1:outer}
For a fixed Q2 label let $(\overline w,\overline z)
=(-\overline U_1,\overline U_2)$ be the zero-noise solution, and
$(w,z)=(-\widehat U_1,\widehat U_2)$ the finite-noise target-only one.
Let $L\ge1$, $T>0$, and
\[
 \epsilon_L(T)=\left[\frac{3q^2}{4}+\frac{K(-L)}{2L}\right]T.
\]
If the normalized zero-noise coordinates satisfy
\[
 \overline c(T)\ge2qL,\qquad \rho\overline s(0)\ge2qL,
 \qquad \epsilon_L(T)\le\frac{\log2}{4},
\]
then the finite-noise label stays in Q2 with $c\ge qL$ and
$\rho s\ge qL$ through $T$, and
\[
 |\log(w/\overline w)|\le\epsilon_L(T),\qquad
 0\le\log(z/\overline z)\le\epsilon_L(T).
\]
In particular the radial log error is at most $2\epsilon_L(T)$,
and the log error of $w/z$ is at most $2\epsilon_L(T)$.
\end{proposition}
\begin{proof}
Until a first exit of $c\ge qL$, $\rho s\ge qL$, formula
\eqref{inc:AN02:q2:v1:coordinates}, with $c=w/(w^2+z^2)^{1/2}$,
$s=z/(w^2+z^2)^{1/2}$, gives
\[
 (\log w)'=d_q-\frac{q}{2c}K(-c/q),\qquad
 (\log z)'=\frac\lambda2+d_q+\frac{\rho q}{2s}K(-\rho s/q).
\]
Here $0\le d_q\le3q^2/4$, and the two extra quotients are bounded
by $K(-L)/(2L)$, since $K$ is increasing and $\lambda<1$.
Integrating proves the displayed log bounds up to a first exit.
The Euclidean radius ratio is between $e^{-\epsilon_L}$ and
$e^{\epsilon_L}$, so each normalized coordinate is at least
$e^{-2\epsilon_L}$ times its zero-noise counterpart.
The latter have $\overline c$ decreasing and $\overline s$ increasing.
At a proposed exit the actual normalized margins therefore exceed
$2qL e^{-(\log2)/2}>qL$. This contradiction closes the interval.
This is a labelwise comparison, not yet a Jacobian or density-error bound.
\end{proof}

\section{Analytical target-only crossing and transit clocks}
\label{inc:AN02:q2:v1:clocks}
All asymptotic statements in this section hold as $q\downarrow0$.
They concern the exact finite-noise target-only flow, not the original
population during strong saturation.

\begin{theorem}[Weak-axis crossing with a renormalized constant]
\label{inc:AN02:q2:v1:cross}
Fix $\rho\in(0,1)$ and $\alpha_0=\pi/2+d_0$, $0<d_0<\pi/2$.
The target-only label crosses $\theta=\pi/2$ at a finite time, and
\begin{equation}\label{inc:AN02:q2:v1:crosstime}
 T_\times(q,\alpha_0)
 =\frac2\lambda\log\frac{\tan d_0}{q}+C_\rho+o(1),
\end{equation}
where the analytically defined finite constant is
\begin{align}
 C_\rho={}&\int_0^1\frac{2\dd z}{K(-z)+\lambda z}\nonumber\\
 &+\int_1^\infty
 \left[\frac{2}{K(-z)+\lambda z}-\frac{2}{\lambda z}\right]\dd z.
 \label{inc:AN02:q2:v1:crossconstant}
\end{align}
The convergence is uniform for $\rho$ and $d_0$ in compact subintervals
of $(0,1)$ and $(0,\pi/2)$, respectively.
\end{theorem}
\begin{proof}
For $d=\theta-\pi/2\in[0,d_0]$,
\begin{align*}
 -v_q(\pi/2+d)
 ={}&\frac\lambda2\sin d\cos d
 +\frac q2\cos d K(-\sin d/q)\\
 &+\frac{\rho q}{2}\sin d K(-\rho\cos d/q)>0.
\end{align*}
The speed is positive and continuous on this compact interval, including
$d=0$, where it is $q\phi(0)/2$. The crossing time is finite.
Set $z=\tan d/q$. Then $z'=-V_q(z)$, with
\begin{align*}
 V_q(z)={}&\frac\lambda2z
 +\frac{\sqrt{1+q^2z^2}}2
       K\left(-\frac{z}{\sqrt{1+q^2z^2}}\right)\\
 &+\frac{\rho qz\sqrt{1+q^2z^2}}2
       K\left(-\frac{\rho}{q\sqrt{1+q^2z^2}}\right).
\end{align*}
Consequently $T_\times=\int_0^{\tan d_0/q}V_q(z)^{-1}\dd z$.
Pointwise $V_q(z)\to[K(-z)+\lambda z]/2$.
For $0\le z\le1$, the positive $K$ term bounds $V_q$ below uniformly
for small $q$, so dominated convergence applies.

For the remaining integral, put $A=(1+\tan^2d_0)^{1/2}$.
On $1\le z\le\tan d_0/q$ write $V_q=\lambda z/2+E_1+E_2$.
The nonnegative corrections obey
\[
 E_1\le\tfrac A2K(-z/A),\qquad
 E_2\le CqzK(-\rho/(qA)),
\]
with $C$ uniformly bounded on the specified parameter compacts. Thus
\[
 0\le\frac2{\lambda z}-\frac1{V_q(z)}
 \le\frac{4(E_1+E_2)}{\lambda^2z^2}.
\]
The $E_1$ majorant is integrable to infinity. The integral of the
$E_2$ majorant is at most
$CqK(-\rho/(qA))\log(\tan d_0/q)\to0$.
Dominated convergence, with the upper-cutoff indicator, proves convergence
of the subtracted integral. Its limit is the second integral in
\eqref{inc:AN02:q2:v1:crossconstant}. The first integral is finite
because $K(0)>0$, and the second is absolutely convergent by the same
bound. Restoring $\int_1^{\tan d_0/q}2/(\lambda z)$ proves the result.
All majorants are uniform on the stated compacts, which also proves
uniformity.
\end{proof}

\begin{theorem}[Transit through Q1]
\label{inc:AN02:q2:v1:transit}
Fix a compact ratio interval $[\rho_-,\rho_+]\subset(0,1)$ and
choose
\[
 R>\frac{4\rho_+\phi(0)}{1-\rho_+^2}.
\]
For sufficiently small $q$, the target-only trajectory after crossing
$e_2$ reaches $\theta=qR$. If $T_R$ is that arrival time measured
from initialization, then
\begin{equation}\label{inc:AN02:q2:v1:strongtime}
 \frac{T_R}{\log(1/q)}\longrightarrow
          \frac2\lambda+\frac4{1-\lambda}.
\end{equation}
For any fixed $\theta_b\in(0,\pi/2)$, its earlier arrival time satisfies
\begin{equation}\label{inc:AN02:q2:v1:fixedtime}
 \frac{T_{\theta_b}}{\log(1/q)}\longrightarrow
          \frac2\lambda+\frac2{1-\lambda}
          =\frac{2}{\lambda(1-\lambda)}.
\end{equation}
Uniformity holds on compact initial Q2 angle intervals, compact ratio
intervals as above, and compact fixed-angle intervals for $\theta_b$.
\end{theorem}
\begin{proof}
In Q1 the exact clockwise speed is
\begin{align*}
 B_q(\theta)=-v_q(\theta)
 ={}&\frac{1-\lambda}{2}\sin\theta\cos\theta
 +\frac q2\sin\theta K(-\cos\theta/q)\\
 &-\frac{\rho q}{2}\cos\theta K(-\rho\sin\theta/q).
\end{align*}
For $\theta\in[qR,\pi/2]$ and sufficiently small $q$,
$\sin\theta\ge qR/2$, so the choice of $R$ implies
\[
 B_q(\theta)\ge\frac{1-\lambda}{4}\sin\theta\cos\theta.
\]
The speed is strictly positive before the axis and equals
$q\phi(0)/2>0$ at the axis itself. It therefore has no zero on
this compact travel interval, proving arrival.

Set $L(q)=\sqrt{\log(1/q)}$. On
$[qL,\pi/2-qL]$, the ratio of $B_q$ to
$B_0=(1-\lambda)\sin\theta\cos\theta/2$ tends uniformly to one.
Indeed the absolute relative correction is at most
\[
 \frac{qK(-\cos\theta/q)}{(1-\lambda)\cos\theta}
 +\frac{\rho qK(-\rho\sin\theta/q)}{(1-\lambda)\sin\theta},
\]
and both projected margins grow proportionally to $L$.
Since $\int\dd\theta/(\sin\theta\cos\theta)=\log\tan\theta$,
the interior travel time is
\[
 \frac4{1-\lambda}\log(1/q)+o(\log(1/q)).
\]
The strong-end layer $[qR,qL]$ has time $O(1+\log L)$ by the
lower speed bound. At the weak end, write $d=\pi/2-\theta$.
For $q\le d\le qL$ the same lower bound gives time $O(\log L)$.
For $0\le d\le q$, one has $\sin\theta\ge1/2$ and
$\cos\theta/q\le1$. The positive $K$ term is at least
$qK(-1)/4$, whereas the negative term has absolute value at most
$q^2\phi(0)/2$. Hence, for small $q$, $B_q\ge qK(-1)/8$ on
this last layer, whose duration is $O(1)$.
These layers are $o(\log(1/q))$, proving a post-crossing travel
coefficient $4/(1-\lambda)$. For a fixed $\theta_b$, only the
weak-end logarithm remains, giving $2/(1-\lambda)$. Add Theorem
\ref{inc:AN02:q2:v1:cross}. Each bound is uniform on the stated compacts.
\end{proof}

\paragraph{What these clocks do not imply.}
These theorems prove logarithmic clock coefficients for the explicitly
defined target-only flow. They do not establish that the original
population follows that flow through saturation, that its Q2 retention
coefficient vanishes, or that no joint parameter regime can satisfy an
initialized seed bound. The original residual changes when the strong
concept is learned. A post-saturation return theorem needs that changed
field; it cannot be inferred from a target-only absorption time.

\section{A localized original-flow comparison on the Q2 side}
\label{inc:AN02:q2:v1:tracking}
We now use the \emph{original} balanced population. At any time choose a
nonnegative locally bounded scalar $\mathfrak M(t)$ and decompose the
full output exactly as
\begin{equation}\label{inc:AN02:q2:v1:decomposition}
 f_t(X)=\mathfrak M(t)X_1e_1+g_t(X),\qquad
 \mathcal E(t)=\sqrt{a_*}\,\|g_t\|_{L^2(P)}.
\end{equation}
The scalar $\mathfrak M$ is an analytical comparator coefficient; it
is not identified with a chart mass without proof. This decomposition
retains every neuron through $g_t$. It makes no statement that
$\mathcal E$ is already small along the initialized trajectory.

\begin{lemma}[Left-facing Gaussian suppression]
\label{inc:AN02:q2:v1:gaussiansuppression}
For $q\le1$, $a_1\le0$, and a unit vector $a$,
\begin{equation}\label{inc:AN02:q2:v1:suppression}
 |(A_\theta a)_1|\le q,\qquad
 |a_1|\,\|A_\theta e_1\|\le2q.
\end{equation}
Consequently, for $C_f$ in Section 1 and unit $b$,
\begin{align}
 \|C_fa\|&\le q\mathfrak M+\mathcal E,\nonumber\\
 \|C_f^Tb\|&\le2q\mathfrak M
           +a_*\mathfrak M\|b-a\|+\mathcal E.
 \label{inc:AN02:q2:v1:Cbound}
\end{align}
\end{lemma}
\begin{proof}
The coordinate formula gives
\[
 (A_\theta a)_1=\frac q2K(a_1/q)
 +\frac{q^2a_1}{2}[\Phi(a_1/q)+\Phi(\rho a_2/q)].
\]
Since $K(a_1/q)\le\phi(0)<1/2$ and the bracket is at most $3/2$,
its absolute value is at most $q\phi(0)/2+3q^2/4\le q$.
For the second bound put $\beta=-a_1\ge0$, and use raw cluster
matrices $A_p=\E_{P_p}[XX^T\mathbf1_{a^TX>0}]$.
The exact Gaussian trace on cluster 1 is
\[
 \operatorname{tr}A_1=(1+2q^2)\Phi(-\beta/q)
                         -\beta q\phi(\beta/q).
\]
Since $\beta\Phi(-\beta/q)\le q\phi(\beta/q)$,
$\beta\|A_1e_1\|\le3q\phi(0)$. On cluster 2, Cauchy--Schwarz gives
$\|A_2e_1\|\le q\sqrt{\lambda+2q^2}\le\sqrt3q$.
Average these bounds and use $\beta\le1$ to obtain
$\beta\|A_\theta e_1\|\le(3\phi(0)+\sqrt3)q/2<2q$.
The case $\beta=0$ is included without dividing by $\beta$.

From \eqref{inc:AN02:q2:v1:decomposition},
$C_f=\mathfrak M e_1e_1^TA_\theta+C_g$, with
$\|C_g\|\le\sqrt{a_*}\|g_t\|_{L^2(P)}=\mathcal E$ by
Cauchy--Schwarz. Apply \eqref{inc:AN02:q2:v1:suppression} to
$C_fa$ and $C_f^Ta$. Replacing $a$ by $b$ in the latter costs at most
$a_*\mathfrak M\|b-a\|$, giving \eqref{inc:AN02:q2:v1:Cbound}.
\end{proof}

\begin{theorem}[Exact conditional Q2-side tracking]
\label{inc:AN02:q2:v1:originaltracking}
Take the target-only reference with the same initial label as the original
characteristic. Define nonnegative scalar envelopes, initially zero, by
\begin{align}
 \Delta'&=a_*\mathfrak M\Delta+3q\mathfrak M+2\mathcal E,
 \label{inc:AN02:q2:v1:envelopes}\\
 H_Q'&=a_*H_Q+a_*(1+\mathfrak M)\Delta
                       +2q\mathfrak M+\mathcal E,\nonumber\\
 L_Q'&=2a_*(\Delta+H_Q)+2q\mathfrak M+2\mathcal E.\nonumber
\end{align}
On any interval on which the original input satisfies $a_1\le0$
and $\Delta<\pi$, continuous angular lifts satisfy
\begin{equation}\label{inc:AN02:q2:v1:trackingbounds}
 |\psi-\theta|\le\Delta,\quad
 |\theta-\widehat\theta|\le H_Q,\quad
 |\psi-\widehat\theta|\le\Delta+H_Q,\quad
 |\log(m/\widehat m)|\le L_Q.
\end{equation}
These bounds include full population feedback through the explicit defect
$\mathcal E$. They are not an unconditional initialized bound on that
defect, nor a post-crossing tracking theorem.
\end{theorem}
\begin{proof}
For $\delta=\psi-\theta$, the exact symmetry identity gives
\[
 \delta'=-\operatorname{tr}(A_\theta)\sin\delta
 -(b^\perp)^TC_fa+b^TC_fa^\perp.
\]
For $|\delta|<\pi$ the target term points inward. Lemma
\ref{inc:AN02:q2:v1:gaussiansuppression} and $\|b-a\|\le|\delta|$
give
\[
 D^+|\delta|\le a_*\mathfrak M|\delta|
                         +3q\mathfrak M+2\mathcal E.
\]
Comparison closes the angular-difference guard wherever $\Delta<\pi$.
Next, the exact input velocity satisfies
\[
 |\theta'-v_q(\theta)|
 \le a_*(1+\mathfrak M)|\delta|+2q\mathfrak M+\mathcal E.
\]
The target identity $A_\theta'a_\theta=0$ gives $|v_q'|\le a_*$.
Comparing with $\widehat\theta'=v_q(\widehat\theta)$ gives the
$H_Q$ bound. Put $g_0(\theta)=a_\theta^TA_\theta a_\theta$,
so $(\log\widehat m)'=2g_0(\widehat\theta)$ and
$|g_0'|\le a_*$. Finally
\[
 |b^T(A_\theta-C_f)a-g_0(\widehat\theta)|
 \le a_*(\Delta+H_Q)+q\mathfrak M+\mathcal E.
\]
Integrating twice this bound yields $L_Q$. The output-angle bound is
triangular. Every estimate holds up to first exit from the left-facing
input region; its persistence is not assumed beyond that stopping time.
\end{proof}

\section{What the retention argument still must prove}
\label{inc:AN02:q2:v1:residue}
The principal strong output has been separated from the error. To turn
Theorem \ref{inc:AN02:q2:v1:originaltracking} into a retained-seed law,
one still needs an initialized, quantitatively useful bound on $g_t$,
with the time-weighted envelopes in \eqref{inc:AN02:q2:v1:envelopes}
under control. The theorem pays for strong angular spread, weak output,
other labels, and any incorrect strong-output approximation through
this exact defect. It does not assert these contributions are small.

For parameter powers placing bulk Q2 labels across the weak axis before
strong learning, a Q2-side comparison alone is insufficient. A positive-side
transit and changed-residual return theorem is then required, at a
correctly matched later chart-entry time. The coordinate energy
$\widehat U_2^2$ is useful for that argument but is not the radial chart
mass $N$ in the leading logistic system until the geometry is controlled.

Likewise, the target-only clock theorems do not prove impossibility of
an initialized global-rate route or possibility of a local-rate route
in every suggested parameter window. Those conclusions require uniform
entry laws, their $q$-powers, tail moments, and finite-noise transfer.
The companion local-passage increment proves a conditional nonlinear
rate and collective-mode estimate; it does not fill those entry gaps.

\begin{thebibliography}{9}\small\raggedright
\setlength{\itemsep}{2pt}\setlength{\parskip}{0pt}
\bibitem{F}\path{RELU_POPULATION_FOUNDATIONS_v2.tex}, supplied exact
balanced characteristic equations and Gaussian integration identities.
\bibitem{P}\path{THEOREM_A_ANALYTICAL_PROGRESS.tex}, supplied checkpoint,
\texttt{pa:earlybounds}, \texttt{pa:targetcomparison}, \texttt{pa:setup}.
\bibitem{A2}\path{AN02_strong_ancestry_and_tail_transport_v1.tex},
finite-noise positive-coordinate ancestry inequality and the separate
zero-noise comparison.
\bibitem{A3}\path{AN03_two_sided_passage_and_Q2_retention_v2.tex},
separately repaired original AN03 increment.
\bibitem{L}\path{AN03_local_rate_mean_tracking_v1.tex}, companion
local-rate and collective-mode passage theorem.
\bibitem{R}\path{Pasted markdown.md}, user's current review. Its
numerical tables and matching conjectures are not proof premises.
\end{thebibliography}
\end{document}
